% Section 6.2, from lectures/L06/appendix/b_thermal_drift.md.

\section{Thermal drift, the emitter resistor, and what to build}
\label{c6:app:b}

The problem the emitter resistor solves, an argument for it that does not work, and the one that does.

\subsection{The problem}
\label{c6:sec:b1}
A transistor's base-emitter voltage falls about 2 mV per degree at a fixed collector current, because the saturation current rises steeply with temperature and drags $V_{BE}$ down.

Turn that around. \textbf{At a fixed base-emitter voltage the collector current rises about 8 per cent per degree,} because $e^{2/26} - 1 = 0.08$.

\bookfigure{lectures/L06/appendix/images/drift_against_temperature.png}{Collector current against temperature rise, on a logarithmic vertical axis, for a stage with no emitter resistor and one with a one kilohm emitter resistor. The undegenerated curve climbs by a factor of ten over thirty degrees; the degenerated one is nearly flat.}{c6:fig:driftagainsttemperature}

\textbf{Over a thirty-degree rise that is a factor of ten.} A stage biased at 1 mA on a bench at 20 degrees is at 10 mA inside an enclosure at 50, sitting hard against one end of its load line.

\subsection{An argument that does not survive its own numbers}
\label{c6:sec:b2}
A tempting way to explain what the emitter resistor does:

\begin{enumerate}
  \item A one degree rise lifts the collector current by about 10 per cent, from 1.06 mA to 1.17.
  \item The emitter voltage rises with it, by a tenth of the 1.06 V it was already dropping, so by about 100 mV.
  \item The base is held by the divider, so the base-emitter voltage falls by that same 100 mV, from 0.65 V to about 0.55.
  \item A lower base-emitter voltage means less current, so the operating point is restored.
\end{enumerate}
Every step follows from the one before it, the conclusion is correct, and the argument is wrong.

\textbf{The 10 per cent is right and the rest does not follow.} 100 mV below 0.65 V, at 60 mV a decade:

\[
  \frac{I_C(0.55)}{I_C(0.65)} = e^{-0.100/0.026} = \frac{1}{47}
\]

\textbf{The collector current would fall by a factor of 47, not by 10 per cent.} The circuit never reaches that state: long before the emitter has risen 100 mV the current has been pulled back. The argument describes a system that overshoots wildly and calls it stable. What is wrong is not the physics but the size of the step: \textbf{it applies an open-loop change to a closed-loop circuit.}

\subsection{What actually happens}
\label{c6:sec:b3}
The emitter resistor is a \textbf{local feedback loop}, and the right way to analyse it is Chapter~\ref{c4:ch:feedback}'s.

The base is held by a stiff divider. The disturbance is $V_{BE}$ drifting down at 2 mV per degree, equivalent to the base being driven 2 mV \emph{up}. The emitter follows within a few millivolts, so the emitter voltage rises 2 mV, and

\[
  \Delta I_C = \frac{2\ \text{mV}}{R_E} = \frac{2\ \text{mV}}{1\ \text{k}\Omega} = 2\ \mu\text{A}
\]

On 0.93 mA that is \textbf{0.21 per cent per degree,} against 8 per cent without the resistor, and the base-emitter voltage moves by about 2 mV rather than 100. \textbf{The suppression factor is the loop gain:}

\[
  1 + \frac{R_E}{r_e} = 1 + \frac{1000}{28} = 37
\]

\bookfigure{lectures/L06/appendix/images/drift_against_re.png}{Drift suppression and the emitter factor plotted against the emitter resistor on logarithmic axes, the two curves lying on top of each other above about fifty ohms, with the 220 mV design point marked at 220 ohms and a factor of ten.}{c6:fig:driftagainstre}

\textbf{That number is the emitter factor,} which Chapter~\ref{c7:ch:smallsignal} introduces as the thing that decides gain. The emitter resistor divides the drift by exactly the factor by which it reduces the gain: \textbf{stability costs gain, one for one.} B.2's story of voltages in sequence has no place to put a loop gain.

\subsection{The 220 millivolt rule}
\label{c6:sec:b4}
Enough emitter resistance for the suppression wanted and no more, since every ohm costs gain. A good rule is to drop about 220 mV across it:

\[
  R_E = \frac{220\ \text{mV}}{I_C}
\]

At 1 mA that is 220 ohm, an E12 value exactly; at 10 mA it is 22 ohm; at 20 mA, 11 ohm. The rule is chosen so that the answer lands on the E12 grid across the decade of currents this book uses.

\[
  EF = \frac{r_e + R_E}{r_e} = \frac{26 + 220}{26} = 9.5
\]

A decade of drift suppression bought with a decade of gain. \textbf{22 mV would give 1.85,} barely any suppression; \textbf{2.2 V would give 85,} and a stage divided by 85 is not an amplifier worth having.

\subsection{What to build}
\label{c6:sec:b5}

\subsubsection{\code{ael/bias/point.hpp}}
\begin{cppcode}
namespace ael::bias
{
struct Point
{
    double baseVoltage{};
    double emitterVoltage{};
    double collectorVoltage{};
    double collectorCurrent{};
};

/// The quiescent point of a divider-biased stage, including the base current's droop.
[[nodiscard]] Point quiescentPoint(double supply, double upper, double lower, double emitter,
                                   double collector, double beta = 50.0, double vbe = 0.65) noexcept;

/// Divider current over base current. Ten is the usual rule of thumb.
[[nodiscard]] double stiffness(double supply, double upper, double lower,
                               double collectorCurrent, double beta = 50.0) noexcept;

[[nodiscard]] double driftWithoutDegeneration() noexcept;
[[nodiscard]] double driftWithDegeneration(double collectorCurrent, double emitterResistor) noexcept;
[[nodiscard]] double driftSuppression(double collectorCurrent, double emitterResistor) noexcept;
[[nodiscard]] double degenerationResistor(double collectorCurrent, double drop = 0.220) noexcept;
}
\end{cppcode}

\code{quiescentPoint} must account for the droop of \secref{c6:sec:a3}. It is self-referential, so solve it: two or three fixed-point iterations converge, and a closed form exists if you prefer.

\textbf{Put the emitter current, not the collector current, through the emitter resistor.} Once the base current is in the calculation it is in both places, so $I_E = (\beta + 1)I_B$ and $emitterVoltage = I_E R_E$. At $\beta = 50$ the difference is two per cent, 19 mV here, and the shipped suite checks it to a tenth of a millivolt.

\subsubsection{Temperature in \code{ael/device/bjt.hpp}}
Add a \code{temperature} in kelvin to \code{Parameters}, defaulting to 300.15, and make both the thermal voltage and the saturation current depend on it:

\[
  V_T = \frac{kT}{q}, \qquad
  I_S(T) = I_S(T_0)\left(\frac{T}{T_0}\right)^3 \exp\left[\frac{E_g q}{k}\left(\frac{1}{T_0}-\frac{1}{T}\right)\right]
\]

with $E_g = 1.11$ eV. \textbf{That is where the $-2$ mV per degree comes from; it is not put in by hand.}

\textbf{The two together give about $-1.77$ mV per degree at 1 mA,} not exactly $-2$. The round figure is what every textbook quotes and what to calculate with; the physics gives a number that depends on the current, and the Cross-check is partly about that difference.

\subsubsection{What good looks like}
About sixty lines. The temperature addition is six.

\subsection{What this section is blind to}
\label{c6:sec:b6}
\begin{itemize}
  \item \textbf{Self-heating.} The transistor's own dissipation raises its temperature, so drift is partly a loop through the package. Negligible here; dominant in Chapter~\ref{c8:ch:follower}'s output stage.
  \item \textbf{The divider's own temperature coefficient.} Resistors drift about 100 parts per million per degree, a hundredth of the transistor, and are ignored.
  \item \textbf{Beta's temperature coefficient.} It rises about 0.5 per cent per degree, moving the droop of \secref{c6:sec:a3}: second order on first.
\end{itemize}
